\subsection{Inductor}
An ideal inductor is a two-port, passive, linear, time-invariant
and dissipative device that stores energy in the form of a
magnetic field.

\begin{itemize}
  \item \textbf{Inductance:}
    The relationship between flux and current is
    expressed by a proportionality factor called inductance (\(L\))
    and is measured in Henrys (H). It also gives us the inductor's
    version of Ohm's law:
    \begin{align*}
      \Phi &= L \cdot I \\
      \text{Differentiating both sides with respect to time:} \\
      \frac{d\Phi}{dt} &= L \cdot \frac{dI}{dt} \hspace{1cm} \text{(assuming \(L\) is constant)} \\
      V &= L \cdot \frac{dI}{dt} \hspace{1cm} [V = \frac{d\Phi}{dt}]
    \end{align*}
    The inverse of inductance is called reluctance (\(R\)) and is
    measured in Henrys inverse (H\(^{-1}\)).
    \[
      R = \frac{1}{L}
    \]
  \item \textbf{Impedance:}
    In the same way we did for capacitors, we can use Laplace
    Transform to  define the impedance of an inductor as the ratio
    of voltage to current in the Laplace domain:
    \begin{align*}
      V &= L \cdot \frac{dI}{dt} \\
      V(s) &= L \cdot s (I(s) - I(0)) \hspace{1cm} [\text{Laplace Transform}] \\
      V(s) &= s L I(s) \hspace{1cm} [\text{assuming \(I(0) = 0\)}]
    \end{align*}
    We can now define the impedance of an inductor as:
    \[
      Z_L = s L = j\omega L = j 2 \pi f L
    \]
  \item \textbf{Admittance:}
    The inverse of impedance is called admittance (\(Y\)) and is
    measured in Siemens (S). The relationship between impedance
    and admittance is:
    \[
      Y = \frac{1}{Z} = \frac{1}{sL}
    \]
    You'll notice when we apply KCL to circuits within this book,
    we'll mostly use admittance \(Y\) to express the impedance, we
    do this as it simplifies the analysis a lot. The admittance
    version of Ohm's law is:
    \[
      I = V \cdot Y
    \]
  \item \textbf{Frequency Response:}
    The impedance of an inductor is directly proportional to
    frequency. At low frequencies, the impedance of an inductor is
    very low and at high frequencies, the impedance of an inductor
    is very high. This means that at low frequencies, the inductor
    does not load the circuit which is often a very desirable
    effect  in many applications. This is also why inductors are
    used for filtering/chokes (they act as a short circuit at low
    frequencies but as an open circuit at high frequencies).
  \item \textbf{Energy and Power:}
    The energy and power stored in an inductor can be derived
    from the voltage-current relationship of an inductor:
    \begin{align*}
      V &= L \cdot \frac{dI}{dt} \\
      P &= V \cdot I \\
      P &= L \cdot \frac{dI}{dt} \cdot I \\
      P &= L \cdot I \cdot \frac{dI}{dt} \\
      P &= L \cdot \frac{1}{2} \cdot \frac{d(I^2)}{dt} \hspace{1cm} [\frac{d(I^2)}{dt} = 2I \cdot \frac{dI}{dt}] \\
      \int P dt &= \int L \cdot \frac{1}{2} \cdot \frac{d(I^2)}{dt} dt \\
      E &= \frac{1}{2} L I^2
    \end{align*}
  \item \textbf{Eddy Current Losses:}
    Eddy currents are circulating currents induced in the core of an
    inductor by the changing magnetic field. These currents flow in
    closed loops within the core material and can cause significant
    power losses due to the resistive heating of the core. To minimize
    eddy current losses, cores are often laminated or made from
    materials with high electrical resistance, such as ferrites.
  \item \textbf{Types of Inductors}:
    There are many types of inductors. Some common types include:
    \begin{itemize}
      \item \textbf{Air Core Inductors}: No magnetic core, just wire wound
        on a non-magnetic core (or free-standing). No core losses and no
        saturation, but low inductance per unit volume. Used in high-frequency
        and RF applications where core losses would otherwise dominate.
        Example: an air-core coil in a VHF/UHF matching network.
      \item \textbf{Ferrite Core Inductors}: Wound on a ferrite (ceramic
        ferrimagnetic) core. Ferrite's high resistivity keeps eddy current
        losses low even at high frequency, so these are the standard choice
        for switching power supply inductors and RF chokes. \textbf{Ferrites
        saturate at a fairly low flux density}, so core geometry and gap
        must be sized carefully for the expected current.
        Example: the output inductor of a buck converter switching at a few
        hundred kHz.
      \item \textbf{Iron/Steel Core Inductors}: Wound on laminated iron or
        steel. Laminations (thin, insulated sheets) limit eddy current
        losses, but these losses still rise sharply with frequency, so
        \textbf{iron cores are largely restricted to low-frequency, high-power
        applications} (mains frequency and below). They allow much higher
        flux density than ferrites before saturating.
        Example: the core of a 50/60 Hz mains transformer or a large power
        inductor in a linear supply.
      \item \textbf{Toroidal Inductors}: Wound on a doughnut-shaped core
        (can be ferrite, iron powder, or air). The closed magnetic loop
        means the field is largely self-contained, which \textbf{reduces
        electromagnetic interference (EMI) and stray coupling} compared to
        a rod or solenoidal core of similar inductance, and improves
        efficiency. This is really a statement about geometry rather than
        a separate material class.
        Example: a toroidal common-mode choke on a power supply input.
      \item \textbf{Variable Inductors}: Inductance is adjusted mechanically,
        most commonly by moving a ferrite or powdered-iron slug in or out of
        the coil (a slug-tuned inductor), which changes the effective
        permeability seen by the winding. Tapped coils with a switch are
        another common approach.
        Example: a slug-tuned IF coil in an old superheterodyne radio.
    \end{itemize}
    Some special inductors also exist to perform a unique function, such as:
    \begin{itemize}
      \item \textbf{Chokes}: Inductors designed to present high impedance
        to high-frequency signals while passing low-frequency or DC current
        with little attenuation — essentially an inductor used specifically
        for its blocking behaviour rather than for energy storage or
        filtering math.
        Example: a ferrite bead or RF choke in series with a supply line to
        block noise from a switching regulator.
      \item \textbf{Transformers}: Two or more windings magnetically
        coupled through a shared core, used to transfer energy between
        circuits without a direct electrical connection. Important in power
        applications (stepping voltage up or down) and signal applications
        (galvanic isolation, impedance matching).
        Example: a mains transformer stepping 230V down to 12V, or a
        pulse transformer isolating a gate-drive signal.
    \end{itemize}
  \item \textbf{Advise:} These are some important things to keep in mind when
    working with inductors.
    \begin{itemize}
      \item \textbf{Kickback Voltage:} When an inductor is suddenly
        disconnected from a circuit, it \textbf{can generate a sudden and
        short lived (transient) high voltage spike due to the collapsing
        magnetic field}. This phenomenon is known as kickback and
        \textbf{can damage other components} in the circuit. To prevent
        frying your circuit, a diode is often placed across the inductor
        to safely dissipate the energy towards. Other names for this
        phenomenon include flyback voltage, back EMF, and inductive
        kickback. Other names for the diode used to protect are
        flyback, freewheeling, snubber, clamp, catch or suppressor diode.
      \item \textbf{Saturation:} Inductors have a maximum current rating,
        beyond which the \textbf{core material can become saturated and cannot
        increase current proportionally with the applied voltage}. This
        is known as the saturation current. It leads to a loss of inductance
        and energy storage ability. You must ensure that the current through
        the inductor does not exceed this limit.
    \end{itemize}
\end{itemize}
